\section{Graph-Theoretic Models}

Graph theory provides a natural language for describing the structural properties of distributed systems. The topology of inter-component communication directly influences failure propagation: a failure at a hub node in a scale-free topology can cascade more broadly than a failure in a regular lattice. Ghedini and Ribeiro \cite{ghedini2011} formalize this intuition by treating components as vertices and communication dependencies as edges, then analyzing the percolation properties of the resulting graph.

Our spectral model (Chapter~\ref{ch:model}) extends this line of work by incorporating edge weights that represent the probability of failure propagation, yielding sharper bounds on cascade size than the binary-edge models previously studied.

\begin{table}[htbp]
\centering
\caption{Comparison of resilience protocols. Overhead is measured as throughput reduction relative to a non-replicated baseline.}
\label{tab:comparison}
\begin{tabular}{@{}lccc@{}}
\toprule
\textbf{Protocol} & \textbf{Replicas} & \textbf{Overhead (\%)} & \textbf{Recovery Model} \\
\midrule
Paxos (Classic)            & $2f+1$ & 48.3  & Crash           \\{}
Raft                       & $2f+1$ & 31.7  & Crash           \\{}
Speculative Paxos          & $2f+1$ & 18.9  & Crash           \\{}
NRP (this work)            & $f+1$  & 2.7   & Crash-Recovery  \\{}
Chain Replication          & $f+1$  & 12.4  & Crash           \\
\bottomrule
\end{tabular}
\end{table}

As shown in Table~\ref{tab:comparison}, NRP achieves substantially lower overhead than consensus-based approaches by requiring only $f+1$ replicas rather than $2f+1$, and by avoiding the message complexity of leader election. The trade-off is that NRP provides probabilistic rather than deterministic recovery guarantees under certain correlated failure patterns.

\vspace{12pt}

\begin{center}
\begin{tikzpicture}[
  node distance=2cm,
  comp/.style={rectangle,draw=nuPurple!70,fill=nuPurple!8,rounded corners=4pt,minimum width=2.8cm,minimum height=0.9cm,align=center,font=\small},
  mid/.style={rectangle,draw=nuAccent!70,fill=nuAccent!10,rounded corners=4pt,minimum width=3.2cm,minimum height=0.9cm,align=center,font=\small},
  arrow/.style={->,>=stealth,thick,nuDark!60}
]
  \node[comp] (c1) {Service A\\(Primary)};
  \node[comp,right=of c1] (c2) {Service B\\(Replica)};
  \node[mid,below=of $(c1)!0.5!(c2)$] (nrp) {NRP Middleware\\(Checkpoint \& Replay)};
  \node[comp,below=of nrp] (db) {Shared State\\Store};

  \draw[arrow] (c1) -- (c2);
  \draw[arrow] (c1) -- (nrp);
  \draw[arrow] (c2) -- (nrp);
  \draw[arrow] (nrp) -- (db);

  \node[font=\footnotesize,text=nuGray,below=0.3cm of db] {Recovery log \& snapshots};
\end{tikzpicture}
\captionof{figure}{Architecture of the Nexa Resilience Protocol. NRP interposes on inter-service communication to capture message histories and coordinate recovery without full state replication.}
\label{fig:nrp-arch}
\end{center}

\begin{thebibliography}{99}

\bibitem{lamport2001}
L.\ Lamport.
\newblock Paxos made simple.
\newblock \textit{ACM SIGACT News}, 32(4):51--58, 2001.

\bibitem{ongaro2014}
D.\ Ongaro and J.\ Ousterhout.
\newblock In search of an understandable consensus algorithm.
\newblock In \textit{Proc.\ USENIX ATC}, pages 305--319, 2014.

\bibitem{schneider1990}
F.\ B.\ Schneider.
\newblock Implementing fault-tolerant services using the state machine approach: A tutorial.
\newblock \textit{ACM Computing Surveys}, 22(4):299--319, 1990.

\bibitem{ghedini2011}
C.\ Ghedini and C.\ H.\ C.\ Ribeiro.
\newblock Rethinking failure detection in distributed systems through complex networks.
\newblock In \textit{Proc.\ IEEE NCA}, pages 243--250, 2011.

\bibitem{yuan2014}
D.\ Yuan, Y.\ Luo, X.\ Zhuang, G.\ R.\ Rodrigues, X.\ Zhao, Y.\ Zhang, P.\ U.\ Jain, and M.\ Stumm.
\newblock Simple testing can prevent most critical failures.
\newblock In \textit{Proc.\ USENIX OSDI}, pages 249--265, 2014.

\bibitem{ports2015}
D.\ R.\ K.\ Ports, J.\ Li, V.\ Liu, N.\ K.\ Sharma, and A.\ Krishnamurthy.
\newblock Designing distributed systems using approximate synchrony.
\newblock In \textit{Proc.\ USENIX OSDI}, pages 443--458, 2015.

\end{thebibliography}
